\icode[magenta]{E-MM-SI-L4-261014-002}
1986年，阿瑟·阿什金发明了第一代光镊，经过30多年的发展，光镊技术也越来越成熟，并被广泛运用于医学及生物学等诸多领域，2018年诺贝尔物理学奖被授予美国的亚瑟·阿斯金、法国的杰哈·莫罗和加拿大的唐娜·斯特里克兰三位物理学家，以表彰他们“在激光物理领域的突破性发明”。为简单起见，我们建立如下模型。如图所示，有两束功率均为\(P\)，波长为\(\lambda\)的平行细激光束\(P_1\)和\(P_2\)，与一焦距为\(f\)的透镜主光轴平行射入凸透镜，后进入一球状透明介质，出射透明介质的光束仍与主光轴平行。透镜主光轴在竖直方向，且与球状透明介质的一条	直径重合，此时透明介质球恰好处于静止状态，已知球状透明介质的质量为\(m\)，当地的重力加速度为\(g\)，真空中光速为\(c\)，以下说法正确的有\paren[]
	
\begin{tikzpicture}[scale=.5,decoration={markings,mark=at position 0.5 with {\arrow{>}}},font=\scriptsize]
	\draw[postaction={decorate},magenta] (-2,2)--(-2,-.3);
	\draw[postaction={decorate},magenta] (2,2)--(2,-.3);
	\draw[fill=white] (0,0) ellipse(3 and .5);
	\draw[postaction={decorate},magenta] (-2,-.39)--(1,-3);
	\draw[postaction={decorate},magenta] (2,-.39)--(-1,-3);
	\draw[fill=white] (0,-4.75) circle(2);
	\draw[postaction={decorate},magenta] (-1,-3)--(-1.5,-6.05);
    \draw[postaction={decorate},magenta] (1,-3)--(1.5,-6.05);
	\draw[postaction={decorate},magenta] (-1.5,-6.05)--(-1.5,-7.5);
	\draw[postaction={decorate},magenta] (1.5,-6.05)--(1.5,-7.5);
	\draw[dashed,cyan] (0,3)--(0,-8);
	\draw[<->] (-2,1.5)--node[fill=white]{\(L\)}(0,1.5);
	\node at(-2,1)[left]{\(P_1\)};
	\node at(2,1)[right]{\(P_2\)};
\end{tikzpicture}
\begin{choices}
	\item 若激光束\(P_1\)的功率大于激光束\(P_2\)的功率,则透明介质球将向左侧偏移
	\item 若将介质球相对于透镜上移一小段距离,则介质球受到激光的作用力将变小
	\item 若将介质球相对于透镜下移一小段距离,则介质球受到激光的作用力将变小
	\item 题干条件下激光功率为\(P = \dfrac{mgc}{1 - \frac{f}{\sqrt{L^2 + f^2}}}\)
\end{choices}
\begin{answer}
    
\end{answer}